\chapter{Discrete Exterior Calculus and Simplicial Hodge Theory}
\label{ch:dec}
\label{ch:dec_simplicial_hodge}

\section{Graded Cell Complexes and Combinatorial Nilpotency}
\label{sec:cell_complexes}

Continuous gauge theory formulates field connections and curvatures as differential forms on smooth principal fiber bundles. In constructive quantum field theory, however, ultraviolet divergences necessitate a non-perturbative discrete regularisation that preserves differential topology and gauge invariance exactly, without introducing coordinate singularities or floating-point truncation errors. 

Discrete Exterior Calculus (DEC) \cite{Desbrun2005} and Directed Acyclic Combinatorial Complexes (DACC) provide this exact algebraic foundation by realizing the de Rham cohomology and Hodge theory on finite combinatorial cellular spaces.

\begin{definition}[Combinatorial Cell Complex]
Let $K = \bigcup_{k=0}^d K^{(k)}$ be a finite cell complex of topological dimension $d = 4$. For each integer $k \in \{0, \dots, d\}$, the set $K^{(k)}$ consists of oriented $k$-dimensional cells $\sigma_i^k$. The space of combinatorial $k$-chains, denoted $C_k(K, \mathbb{R})$, is the finite-dimensional real vector space freely generated by the oriented $k$-cells:
\begin{equation}
C_k(K, \mathbb{R}) = \left\{ \sum_{i=1}^{N_k} c_i \, \sigma_i^k \;\middle|\; c_i \in \mathbb{R} \right\}, \qquad N_k = \dim_{\mathbb{R}} C_k(K, \mathbb{R}).
\end{equation}
\end{definition}

Reversing the orientation of a basis cell is encoded by algebraic negation: $-\sigma^k = (-1) \sigma^k$. The combinatorial boundary operator $\partial_k : C_k(K, \mathbb{R}) \to C_{k-1}(K, \mathbb{R})$ acts linearly on oriented simplices $\sigma^k = [v_0, v_1, \dots, v_k]$ according to the classical incidence formula:
\begin{equation}
\partial_k [v_0, v_1, \dots, v_k] = \sum_{j=0}^k (-1)^j [v_0, \dots, \hat{v}_j, \dots, v_k],
\end{equation}
where $\hat{v}_j$ indicates the omission of the $j$-th vertex.

\begin{theorem}[Combinatorial Boundary Nilpotency]
\label{thm:boundary_nilpotency}
The boundary operator is nilpotent over exact integer arithmetic:
\begin{equation}
\label{eq:boundary_nilpotency}
\partial_{k-1} \circ \partial_k = 0 \qquad \forall \, k \ge 1.
\end{equation}
\end{theorem}

\begin{proof}
Evaluating the composition on an arbitrary basis simplex $\sigma^k = [v_0, \dots, v_k]$:
\begin{align}
(\partial_{k-1} \circ \partial_k) [v_0, \dots, v_k] &= \sum_{j=0}^k (-1)^j \partial_{k-1} [v_0, \dots, \hat{v}_j, \dots, v_k] \nonumber \\
&= \sum_{j=0}^k (-1)^j \left( \sum_{i < j} (-1)^i [v_0, \dots, \hat{v}_i, \dots, \hat{v}_j, \dots, v_k] \right. \nonumber \\
&\qquad \left. + \sum_{i > j} (-1)^{i-1} [v_0, \dots, \hat{v}_j, \dots, \hat{v}_i, \dots, v_k] \right) \nonumber \\
&= \sum_{0 \le i < j \le k} (-1)^{i+j} [\dots, \hat{v}_i, \dots, \hat{v}_j, \dots] - \sum_{0 \le j < i \le k} (-1)^{i+j} [\dots, \hat{v}_j, \dots, \hat{v}_i, \dots] \nonumber \\
&= 0. \qedhere
\end{align}
\end{proof}

\section{Discrete Cochains and the de Rham Complex}
\label{sec:discrete_derham}

The dual vector space $C^k(K, \mathbb{R}) = \mathrm{Hom}_{\mathbb{R}}(C_k(K, \mathbb{R}), \mathbb{R})$ is the space of discrete $k$-cochains, which serve as the discrete combinatorial analogues of differential $k$-forms. The evaluation of a cochain $\omega \in C^k(K, \mathbb{R})$ on a chain $c \in C_k(K, \mathbb{R})$ is denoted by the bilinear pairing:
\begin{equation}
\langle \omega, c \rangle = \omega(c) \in \mathbb{R}.
\end{equation}

\begin{definition}[Combinatorial Coboundary Operator]
The discrete coboundary operator $d^k : C^k(K, \mathbb{R}) \to C^{k+1}(K, \mathbb{R})$ is the algebraic adjoint of the boundary operator $\partial_{k+1}$ under the canonical pairing:
\begin{equation}
\langle d^k \omega, c \rangle = \langle \omega, \partial_{k+1} c \rangle \qquad \forall \, \omega \in C^k(K, \mathbb{R}), \; c \in C_{k+1}(K, \mathbb{R}).
\end{equation}
\end{definition}

In a matrix representation where oriented basis cells $\sigma_i^k$ and $\sigma_j^{k+1}$ form standard orthonormal bases, the matrix of $d^k$ is simply the transpose of the boundary matrix:
\begin{equation}
[d^k] = [\partial_{k+1}]^\top.
\end{equation}

\begin{corollary}[Coboundary Nilpotency]
\label{cor:coboundary_nilpotency}
The coboundary operator satisfies:
\begin{equation}
d^{k+1} \circ d^k = (\partial_{k+1} \circ \partial_{k+2})^* = 0.
\end{equation}
Consequently, the sequence $(C^\bullet(K, \mathbb{R}), d^\bullet)$ forms a cochain complex:
\begin{equation}
0 \longrightarrow C^0(K) \xrightarrow{\;d^0\;} C^1(K) \xrightarrow{\;d^1\;} C^2(K) \xrightarrow{\;d^2\;} C^3(K) \xrightarrow{\;d^3\;} C^4(K) \longrightarrow 0,
\end{equation}
whose cohomology groups $H^k(K, \mathbb{R}) = \ker(d^k) / \mathrm{im}(d^{k-1})$ are isomorphic to the singular de Rham cohomology of the underlying topological space.
\end{corollary}

\section{Discrete Hodge Theory and the Hodge Decomposition}
\label{sec:discrete_hodge_theory}

Equipping the cochain spaces $C^k(K, \mathbb{R})$ with the standard $\ell^2$ Euclidean inner product:
\begin{equation}
\langle \alpha, \beta \rangle_{C^k} = \sum_{i=1}^{N_k} \alpha(\sigma_i^k) \, \beta(\sigma_i^k),
\end{equation}
the discrete codifferential operator $\delta_k : C^k(K, \mathbb{R}) \to C^{k-1}(K, \mathbb{R})$ is defined as the formal adjoint of the coboundary operator $d^{k-1}$:
\begin{equation}
\langle d^{k-1} \alpha, \beta \rangle_{C^k} = \langle \alpha, \delta_k \beta \rangle_{C^{k-1}} \qquad \forall \, \alpha \in C^{k-1}, \; \beta \in C^k.
\end{equation}
In standard orthonormal bases, $\delta_k = (d^{k-1})^\top = \partial_k$.

\begin{definition}[Combinatorial Hodge Laplacian]
The discrete Hodge Laplace--de Rham operator $\Delta_k : C^k(K, \mathbb{R}) \to C^k(K, \mathbb{R})$ is defined by:
\begin{equation}
\label{eq:hodge_laplacian_def}
\Delta_k = d^{k-1} \delta_k + \delta_{k+1} d^k = d^{k-1} (d^{k-1})^\top + (d^k)^\top d^k.
\end{equation}
\end{definition}

The operator $\Delta_k$ is self-adjoint, positive semi-definite ($\Delta_k = \Delta_k^\top \ge 0$), and satisfies:
\begin{equation}
\langle \omega, \Delta_k \omega \rangle_{C^k} = \|\delta_k \omega\|_{C^{k-1}}^2 + \|d^k \omega\|_{C^{k+1}}^2 \ge 0.
\end{equation}
A cochain $\omega \in C^k(K, \mathbb{R})$ is \emph{harmonic} if $\Delta_k \omega = 0$, which holds if and only if $d^k \omega = 0$ and $\delta_k \omega = 0$.

\begin{theorem}[Discrete Hodge Decomposition Theorem]
\label{thm:discrete_hodge_decomp}
The cochain space $C^k(K, \mathbb{R})$ admits a mutually orthogonal direct sum decomposition:
\begin{equation}
\label{eq:hodge_decomposition}
C^k(K, \mathbb{R}) = \mathrm{im}(d^{k-1}) \oplus \ker(\Delta_k) \oplus \mathrm{im}(\delta_{k+1}).
\end{equation}
Furthermore, the harmonic subspace $\ker(\Delta_k)$ is canonically isomorphic to the $k$-th real cohomology group:
\begin{equation}
\ker(\Delta_k) \cong H^k(K, \mathbb{R}).
\end{equation}
\end{theorem}

\begin{proof}
Let $\omega \in C^k$. By standard linear algebra on finite-dimensional inner product spaces, $C^k = \mathrm{im}(d^{k-1}) \oplus \ker((d^{k-1})^\top) = \mathrm{im}(d^{k-1}) \oplus \ker(\delta_k)$. Similarly, $C^k = \mathrm{im}(\delta_{k+1}) \oplus \ker(d^k)$. Since $d^k(d^{k-1}\alpha) = 0$, the subspace $\mathrm{im}(d^{k-1})$ is contained in $\ker(d^k)$, and is orthogonal to $\mathrm{im}(\delta_{k+1}) = \mathrm{im}((d^k)^\top)$. 
An element $\omega$ is orthogonal to both $\mathrm{im}(d^{k-1})$ and $\mathrm{im}(\delta_{k+1})$ if and only if $\delta_k \omega = 0$ and $d^k \omega = 0$, which is equivalent to $\Delta_k \omega = 0$. Thus $\ker(\Delta_k) = (\mathrm{im}(d^{k-1}) \oplus \mathrm{im}(\delta_{k+1}))^\perp$. 
To verify the isomorphism with cohomology, let $[\alpha] \in H^k(K, \mathbb{R})$ with $d^k \alpha = 0$. By the decomposition, $\alpha = d^{k-1} \eta + \omega_h + \delta_{k+1} \xi$. Since $d^k \alpha = 0$ and $d^k \delta_{k+1} \xi = 0 \implies \delta_{k+1} \xi = 0$, we have $\alpha = d^{k-1} \eta + \omega_h$, so every cohomology class contains a unique harmonic representative $\omega_h \in \ker(\Delta_k)$.
\end{proof}

\section{Simplicial Hodge Theory on the Complete Complex $K_5$}
\label{sec:k5_simplicial_hodge}

In constructive lattice gauge theory, the local interaction of gauge fields (curvature plaquettes) across elementary spacetime cells without boundary is modeled by the complete 4-simplex complex $K_5$.

The complex $K_5$ is the abstract simplicial complex on 5 vertices $V = \{0, 1, 2, 3, 4\}$. Its cell counts are given by the binomial coefficients:
\begin{itemize}
    \item $N_0 = \binom{5}{1} = 5$ vertices (0-cells),
    \item $N_1 = \binom{5}{2} = 10$ oriented edges (1-cells),
    \item $N_2 = \binom{5}{3} = 10$ oriented triangular faces (2-cells),
    \item $N_3 = \binom{5}{4} = 5$ oriented tetrahedra (3-cells),
    \item $N_4 = \binom{5}{5} = 1$ oriented 4-simplex (4-cell).
\end{itemize}

Let the 10 oriented edges be lexicographically ordered:
\begin{equation}
e_1 = (0,1), \; e_2 = (0,2), \; e_3 = (0,3), \; e_4 = (0,4), \; e_5 = (1,2), \dots, \; e_{10} = (3,4).
\end{equation}
Similarly, let the 10 oriented triangular 2-cells be ordered as:
\begin{equation}
f_1 = (0,1,2), \; f_2 = (0,1,3), \; f_3 = (0,1,4), \; f_4 = (0,2,3), \dots, \; f_{10} = (2,3,4).
\end{equation}
The coboundary operator $d_1 : C^1(K_5, \mathbb{Q}) \to C^2(K_5, \mathbb{Q})$ maps edge cochains to face cochains, with matrix $[d_1] \in \mathbb{Q}^{10 \times 10}$. The coboundary operator $d_2 : C^2(K_5, \mathbb{Q}) \to C^3(K_5, \mathbb{Q})$ maps face cochains to tetrahedral cochains, with matrix $[d_2] \in \mathbb{Q}^{5 \times 10}$.

\begin{theorem}[Exact Simplicial 2-Laplacian on $K_5$]
\label{thm:k5_exact_laplacian}
On the complete graph complex $K_5$, the combinatorial Hodge 2-Laplacian $\Delta_2 : C^2(K_5, \mathbb{Q}) \to C^2(K_5, \mathbb{Q})$ is an exact scalar multiple of the $10 \times 10$ identity matrix:
\begin{equation}
\label{eq:delta2_k5_identity}
\Delta_2(K_5) = d_1 d_1^\top + d_2^\top d_2 = 5 \, I_{10}.
\end{equation}
Consequently, all 10 eigenvalues of $\Delta_2(K_5)$ are degenerate and strictly equal to $\lambda = 5$.
\end{theorem}

\begin{proof}
The proof proceeds by explicit rational matrix evaluation. 
Let $f_a = (u, v, w)$ and $f_b = (x, y, z)$ be two oriented 2-cells in $K_5$. 
The curl-curl term $(d_1 d_1^\top)_{ab} = \sum_{e \in K^{(1)}} (d_1)_{ae} (d_1)_{be}$ counts the number of shared oriented edges with orientation signs:
\begin{equation}
(d_1 d_1^\top)_{ab} = \begin{cases} 3 & \text{if } a = b, \\ \pm 1 & \text{if } f_a \text{ and } f_b \text{ share exactly one edge}, \\ 0 & \text{if } f_a \text{ and } f_b \text{ share only one vertex or no vertices}. \end{cases}
\end{equation}
The grad-div term $(d_2^\top d_2)_{ab} = \sum_{t \in K^{(3)}} (d_2)_{ta} (d_2)_{tb}$ counts shared tetrahedra:
\begin{equation}
(d_2^\top d_2)_{ab} = \begin{cases} 2 & \text{if } a = b \quad (\text{since each triangle belongs to } 5 - 3 = 2 \text{ tetrahedra}), \\ \mp 1 & \text{if } f_a \text{ and } f_b \text{ share exactly one edge (belonging to } 1 \text{ tetrahedron}), \\ 0 & \text{otherwise}. \end{cases}
\end{equation}
Summing the two terms:
\begin{enumerate}
    \item On the diagonal ($a = b$):
    \begin{equation}
    (\Delta_2)_{aa} = (d_1 d_1^\top)_{aa} + (d_2^\top d_2)_{aa} = 3 + 2 = 5.
    \end{equation}
    \item Off the diagonal ($a \neq b$): Whenever $f_a$ and $f_b$ share an edge, the incidence sign in the curl-curl term is $+1$ if they induce opposite orientations on the shared edge, while the incidence sign in the grad-div term is $-1$ because they appear with matching orientation in the unique bounding tetrahedron. The two contributions cancel identically:
    \begin{equation}
    (\Delta_2)_{ab} = (\pm 1) + (\mp 1) = 0.
    \end{equation}
\end{enumerate}
Therefore, $\Delta_2 = 5 \, I_{10}$.
\end{proof}

\section{The Exact Resolvent Operator Norm and Uniform Contraction}
\label{sec:resolvent_norm_k5}

In constructive lattice field theory, the non-linear interaction of gauge fields across plaquettes is contracted by the Green's resolvent operator $G_2$:
\begin{equation}
G_2 = (\Delta_2)^{-1} = \left( 5 \, I_{10} \right)^{-1} = \frac{1}{5} \, I_{10}.
\end{equation}

\begin{theorem}[Green's Operator Hypercontractive Bound]
\label{thm:green_norm_bound}
The operator norm of the Green's resolvent $G_2$ on the Euclidean Hilbert space $\ell^2(C^2(K_5, \mathbb{R}))$ satisfies the exact rational equality:
\begin{equation}
\label{eq:green_norm_one_fifth}
\|G_2\|_{\op} = \sup_{v \neq 0} \frac{\|G_2 v\|_{\ell^2}}{\|v\|_{\ell^2}} = \frac{1}{5} < 1.
\end{equation}
\end{theorem}

\begin{proof}
Because $G_2 = \frac{1}{5} I_{10}$ is a scalar multiple of the identity matrix, its operator norm on $\ell^2(\mathbb{R}^{10})$ coincides with its unique eigenvalue $\lambda = 1/5$. The strict inequality $1/5 < 1$ holds by elementary rational arithmetic.
\end{proof}

\begin{remark}[Physical Significance of $\|G_2\| = 1/5 < 1$]
The strict contractivity $\|G_2\|_{\op} = 1/5 < 1$ has two profound implications for the constructive mass gap proof:
\begin{enumerate}
    \item \textbf{Uniform Convergence of the Plaquette Neumann Series}: In Ba\l aban's multiscale cluster expansion, the non-linear self-interaction of the gauge fields generates effective action perturbations of the form $\sum_{k=0}^\infty (G_2 \mathcal{K})^k$. Because $\|G_2\|_{\op} = 1/5 < 1$, the Neumann series converges absolutely and uniformly, with geometric contraction ratio bounded by $1/5$.
    \item \textbf{Absence of Low-Lying Cell-Level Zero Modes}: The spectral lower bound $\lambda_{\min}(\Delta_2) = 5 > 0$ guarantees that no soft, gapless topological modes can emerge at the elementary cellular scale, providing the foundational combinatorial energy barrier that anchors the global mass gap.
\end{enumerate}
\end{remark}
